% Chapter 1, section: pose-and-screws
\section{3-D pose and screw motion}\label{sec:rot:pose-and-screws}

\subsection{From orientation to pose}\label{sec:rot:orientation-pose}
A rigid body has 6 degrees of freedom: 3 for where its frame's origin is, 3 for how the frame is turned. The pose of $\{B\}$ in $\{A\}$ is the pair $(R_{AB},t_{AB})$, where $t_{AB}$ is the origin of $B$ written in $A$-coordinates.

\begin{deriv}{how a point's coordinates change between frames}
The arrow from $A$'s origin to a point $p$ equals (arrow from $A$'s origin to $B$'s origin) $+$ (arrow from $B$'s origin to $p$). In $A$-coordinates the first is $t_{AB}$; the second is known in $B$-coordinates and converted with $R_{AB}$:
\[
\pre{A}{p}=R_{AB}\,\pre{B}{p}+t_{AB}.
\]
Directions (differences of two points) lose the $t_{AB}$: $\pre{A}{d}=R_{AB}\,\pre{B}{d}$.
\end{deriv}

\subsection{Homogeneous coordinates and the transformation matrix}\label{sec:rot:homog}\label{sec:rot:homogeneous-coordinates-transformation}
The map above is affine, not linear, so it does not compose by plain matrix products. Append one coordinate: $1$ for points, $0$ for directions. Then
\[
\bmat{\pre{A}{p}\\1}=\underbrace{\bmat{R_{AB}&t_{AB}\\0&1}}_{T_{AB}}\bmat{\pre{B}{p}\\1},\qquad
\bmat{\pre{A}{d}\\0}=T_{AB}\bmat{\pre{B}{d}\\0}.
\]
The trailing $0$ makes directions immune to translation automatically, which is the point/vector distinction of \S\ref{sec:rot:vectors-are-arrows} built into the arithmetic. The set
\[
\SE=\Big\{T=\bmat{R&t\\0&1}:\ R\in\SO,\ t\in\RR^3\Big\}
\]
is a group (the special Euclidean group): every rigid motion preserves distances and handedness, and every such map of space has this form.

\begin{deriv}{composition and inverse}
Multiplying blocks: $T_{AB}T_{BC}=\bmat{R_{AB}R_{BC}&R_{AB}t_{BC}+t_{AB}\\0&1}=T_{AC}$. Read the translation part: ``$C$'s origin in $B$, rotated into $A$, plus $B$'s origin in $A$''. For the inverse, solve $\pre{A}{p}=R\pre{B}{p}+t$ for $\pre{B}{p}$:
\[
T_{AB}^{-1}=T_{BA}=\bmat{R\T&-R\T t\\0&1}.
\]
The term $-R\T t$ is $A$'s origin written in $B$. A common bug is to write $-t$ instead of $-R\T t$, which is correct only when $R=I$.
\end{deriv}

\textbf{Frame-notation rules.} (1) $T_{AB}$ maps $B$-coordinates to $A$-coordinates \emph{and} is the pose of $B$ in $A$. (2) Adjacent indices must match and cancel: $T_{WC}=T_{WB}T_{BC}$; a product like $T_{WB}T_{CB}$ is a type error. (3) Relative pose between two keyframes: $T_{12}=T_{W1}^{-1}T_{W2}$. (4) If a motion $\Delta$ is known in frame $C$ but you need it in frame $B$, conjugate: $\Delta_B=T_{BC}\Delta_CT_{BC}^{-1}$. Hand--eye calibration $AX=XB$ is exactly this equation with the unknown mount $X=T_{BC}$.

\begin{figure}[h]
\centering
\begin{minipage}{0.56\linewidth}\includegraphics[width=\linewidth]{chain.pdf}\end{minipage}\hfill
\begin{minipage}{0.42\linewidth}
\caption{A chain of frames: the robot base $\{B\}$ in the world, a camera $\{C\}$ on the robot (pitched down $20^\circ$), and an object seen by the camera. Composing the transforms along the chain maps the camera measurement into world coordinates.}
\label{fig:rot:chain}
\end{minipage}
\end{figure}

\begin{example}{camera detection to world map, and back to pixels}
\emph{Given.} Robot at $(2.0,1.0,0)$\,m, yaw $30^\circ$: $T_{WB}=\bmat{R_z(30^\circ)&(2,1,0)\T\\0&1}$. Camera at $(0.25,0,0.6)$ in the base frame, optical convention, pitched down $20^\circ$: $R_{BC}=R_y(20^\circ)\,R_{BC}^{0}$ with $R_{BC}^0$ from the example in \S\ref{sec:rot:orientation}, giving
\[
R_{BC}=\bmat{0&-0.342&0.940\\-1&0&0\\0&-0.940&-0.342},\qquad T_{WC}=T_{WB}T_{BC}=\bmat{0.500&-0.296&0.814&2.217\\-0.866&-0.171&0.470&1.125\\0&-0.940&-0.342&0.600\\0&0&0&1}.
\]
\emph{Forward.} A detection at $\pre{C}{p}=(0.3,-0.1,1.5)$ gives $\pre{B}{p}=T_{BC}\pre{C}{p}=(1.694,-0.300,0.181)$ and $\pre{W}{p}=(3.617,\,1.587,\,0.181)$: 18\,cm above the floor, as it should be for an object seen slightly above the optical axis by a down-tilted camera.

\emph{Backward (map landmark $\to$ pixel).} $T_{CW}=T_{WC}^{-1}$ (use the $R\T,-R\T t$ formula). A mapped landmark $\pre{W}{l}=(4.0,2.5,0.3)$ is at $\pre{C}{l}=(-0.299,-0.482,2.200)$. With pinhole intrinsics $f=600$\,px, $(c_u,c_v)=(320,240)$: $u=320+600\cdot(-0.299/2.2)=238.4$, $v=240+600\cdot(-0.482/2.2)=108.7$. This ``transform, then project'' is the forward model inside every visual localisation and bundle-adjustment residual.
\end{example}

\subsection{Estimating a pose from point correspondences (Kabsch / Umeyama)}\label{sec:rot:estimating-pose-point}
\begin{deriv}{closed-form least-squares rigid alignment}
Given model points $p_i$ and measured points $q_i\approx Rp_i+t$, minimise $\sum_i\|q_i-Rp_i-t\|^2$. Setting the derivative w.r.t.\ $t$ to zero: $t=\bar q-R\bar p$ (centroids must align). With centred points $p_i'=p_i-\bar p$, $q_i'=q_i-\bar q$, the cost becomes $\text{const}-2\sum_i(q_i')\T Rp_i'=\text{const}-2\tr(RH)$ with $H=\sum_ip_i'(q_i')\T$. Take the SVD $H=USV\T$; then $\tr(RH)=\tr\big((V\T RU)S\big)$ and $V\T RU$ is orthogonal, so the trace is largest when $V\T RU=I$:
\[
R=V\,\mathrm{diag}\big(1,1,\det(VU\T)\big)\,U\T,\qquad t=\bar q-R\bar p.
\]
The $\det$ term prevents returning a reflection when the data are noisy or nearly planar.
\end{deriv}
\begin{example}{object pose from five 3-D points}
A 20\,cm object has 5 known model points; a depth camera measures them with 2\,mm noise. Kabsch recovers the pose with rotation error $1.30^\circ$ and translation error $2.2$\,mm (RMS residual 3.8\,mm). Note the large rotation error: angular accuracy is roughly noise$/$size, so for small objects orientation is the weak part of the estimate. The same routine is the inner step of ICP, aligns estimated and ground-truth trajectories before computing ATE, and (with a scale factor, Umeyama) aligns monocular SLAM trajectories.
\end{example}

\subsection{Chasles' theorem: every rigid displacement is a screw}\label{sec:rot:chasles-theorem-every}
\begin{deriv}{rotation about some axis $+$ translation along the same axis}
Take $(R,t)$ with $R$ a rotation by $\theta\neq0$ about $\nh$. Split $t=t_\parallel+t_\perp$ along and across $\nh$. A screw motion (rotate by $\theta$ about the line through point $q$ with direction $\nh$, then slide $d$ along it) maps $p\mapsto R(p-q)+q+d\nh=Rp+(I-R)q+d\nh$. Match with $Rp+t$:
\[
(I-R)q+d\nh=t .
\]
Dotting with $\nh$ (note $\nh\T(I-R)=0$ because $R\T\nh=\nh$) gives $d=\nh\cdot t$. The remaining equation $(I-R)q=t_\perp$ is solvable because on the plane perpendicular to $\nh$, $R$ is a rotation by $\theta\neq0$ and has no fixed vector, so $I-R$ is invertible there. Hence \emph{any} rigid displacement is a screw: rotation $\theta$ about a unique axis, plus translation $d$ along that same axis. Its \textbf{pitch} is $h=d/\theta$ (metres per radian). Pure rotation: $h=0$. Pure translation: $\theta=0$, the limiting case $h=\infty$.
\end{deriv}

\subsection{Twists: the velocity of a rigid body}\label{sec:rot:twists-velocity-rigid}
Differentiate a moving pose $T(t)$, as we did for $R(t)$:
\[
T^{-1}\dot T=\bmat{R\T\dot R&R\T\dot t\\0&0}=\bmat{\skw{\omega_b}&v_b\\0&0},\qquad
\dot TT^{-1}=\bmat{\dot RR\T&\dot t-\dot RR\T t\\0&0}=\bmat{\skw{\omega_s}&v_s\\0&0}.
\]
Both have the form $[\xi]=\bmat{\skw{\omega}&v\\0&0}\in\se$, the Lie algebra of $\SE$, identified with the \textbf{twist} $\xi=(v,\omega)\in\RR^6$.
\begin{itemize}
\item \textbf{Body twist} $\xi_b$: $\omega_b$ and $v_b=R\T\dot t$, the velocity of the body origin, both in body coordinates. This is what wheel odometry and a VIO front-end produce.
\item \textbf{Spatial twist} $\xi_s$: $\omega_s$ and $v_s=\dot t-\omega_s\times t$. $v_s$ is the velocity of the (imaginary) body point that is currently passing through the world origin. Strange, but it makes $\xi_s$ independent of where the body frame is attached.
\end{itemize}
\textbf{Screw axis $\to$ twist.} A unit-speed screw with direction $\nh$, a point $q$ on the axis and pitch $h$ has $\omega=\nh$ and $v=-\nh\times q+h\nh$. The term $-\nh\times q=\nh\times(0-q)$ is the velocity of the origin caused by spinning about an axis through $q$; $h\nh$ is the slide along the axis. For a revolute joint $h=0$; for a prismatic joint $\omega=0$, $v=\nh$.

\subsection{The exponential of a twist}\label{sec:rot:exponential-twist}
\begin{deriv}{$e^{[\xi]\theta}$ generates the screw motion}
Let $\|\omega\|=1$. Powers of the $4\times4$ matrix: $[\xi]^k=\bmat{\skw{\omega}^k&\skw{\omega}^{k-1}v\\0&0}$ for $k\ge1$. Summing the series block-wise:
\[
e^{[\xi]\theta}=\bmat{e^{\skw{\omega}\theta}&G(\theta)v\\0&1},\qquad
G(\theta)=\sum_{k\ge1}\frac{\theta^k}{k!}\skw{\omega}^{k-1}=I\theta+(1-\cos\theta)\skw{\omega}+(\theta-\sin\theta)\skw{\omega}^2,
\]
where \eqref{eq:rot:skewpow} collapses the series exactly as for Rodrigues. Now insert the screw twist $v=-\omega\times q+h\omega$:
\begin{itemize}
\item $G(\theta)(h\omega)=h\theta\,\omega$, because $\skw{\omega}\omega=0$: the slide along the axis is $h\theta$.
\item $G(\theta)(-\skw{\omega}q)=-\theta\skw{\omega}q-(1-\cos\theta)\skw{\omega}^2q-(\theta-\sin\theta)\skw{\omega}^3q=-\sin\theta\skw{\omega}q-(1-\cos\theta)\skw{\omega}^2q=(I-e^{\skw{\omega}\theta})q.$
\end{itemize}
So $e^{[\xi]\theta}=\bmat{R&(I-R)q+h\theta\,\omega\\0&1}$, precisely the screw of Chasles' theorem: rotate about the line through $q$, slide $h\theta$ along it. Special cases: $h=0$ (pure rotation about an off-origin axis, translation $(I-R)q$); $\omega=0,\ \|v\|=1$ (pure translation, $e^{[\xi]\theta}=\bmat{I&v\theta\\0&1}$).
\end{deriv}

For a general rotation vector $\varphi=\omega\theta$ and $\rho=v\theta$, the same result reads $\Exp(\rho,\varphi)=\bmat{\Exp(\varphi)&V(\varphi)\rho\\0&1}$ with $V=I+\frac{1-\cos\|\varphi\|}{\|\varphi\|^2}\skw{\varphi}+\frac{\|\varphi\|-\sin\|\varphi\|}{\|\varphi\|^3}\skw{\varphi}^2$. Note the coupling matrix $V$: \emph{the translation of $\Exp$ is not $\rho$}. Exponentiating rotation and translation separately (a common shortcut) gives a different, frame-dependent motion. \textbf{Log map:} $\theta,\omega$ from $\Log(R)$, then $v=G(\theta)^{-1}t$ with
$G^{-1}=\frac1\theta I-\frac12\skw{\omega}+\big(\frac1\theta-\frac12\cot\frac\theta2\big)\skw{\omega}^2$ (or $v=t/\|t\|,\ \theta=\|t\|$ if $R=I$).

\begin{example}{screw parameters of a gripper motion, and how to interpolate it}
A gripper moves from the identity pose to $T=\bmat{R_z(90^\circ)&(1,1,0.5)\T\\0&1}$. By hand: the axis is along $z$; $d=\nh\cdot t=0.5$; solve $(I-R)q=(1,1,0)$ with $I-R_z(90^\circ)=\bmat{1&1\\-1&1}$ on the $xy$-block, giving $q=(0,1)$. So the motion is a $90^\circ$ turn about the vertical line through $(0,1)$ plus a $0.5$\,m rise along it: pitch $h=0.5/(\pi/2)=0.318$\,m/rad. The $\Log$ formula agrees: $\xi\theta=(1.571,\,0,\,0.5;\ 0,0,1.571)$, i.e.\ $\theta=\pi/2$, $\omega=e_z$, $v=(1,0,0.318)$, $\omega\cdot v=0.318=h$, $\omega\times v=(0,1,0)=q$.

\emph{Interpolation.} $T(s)=\exp(s\,[\xi\theta])$ moves at constant twist along a helix. At $s=0.5$: $R=R_z(45^\circ)$, $t=(0.707,0.293,0.25)$, which is the origin rotated $45^\circ$ about $(0,1)$ and lifted $0.25$. \figref{fig:rot:screw} compares this with the usual ``lerp translation + slerp rotation''. Both are valid paths with the same endpoints, but the screw path does not depend on which point of the gripper you call the origin: changing the tool frame by $X$ gives exactly $T(s)X$. The decoupled path changes shape when the tool frame is moved, because it moves the chosen origin in a straight line and swings every other point around it.
\end{example}
\begin{figure}[h]
\centering\includegraphics[width=\linewidth]{screw.pdf}
\caption{Interpolating the same displacement two ways. (a) Constant twist: every point follows a helix about one screw axis (dashed). (b) Decoupled: the frame origin moves on a straight line while the orientation slerps, so the gripper tip (orange) takes a different, origin-dependent path.}
\label{fig:rot:screw}
\end{figure}

\subsection{Changing the frame of a twist: the adjoint}\label{sec:rot:changing-frame-twist}
\begin{deriv}{$\xi_s=\Ad_T\,\xi_b$}
From the two definitions, $[\xi_s]=\dot TT^{-1}=T(T^{-1}\dot T)T^{-1}=T[\xi_b]T^{-1}$. Multiply out with \eqref{eq:rot:conj}:
\[
T[\xi]T^{-1}=\bmat{R\skw{\omega}R\T&-R\skw{\omega}R\T t+Rv\\0&0}=\bmat{\skw{R\omega}&Rv+t\times R\omega\\0&0}
\ \Rightarrow\
\Ad_T=\bmat{R&\skw{t}R\\0&R}\ \ (\text{ordering }(v,\omega)).
\]
The angular part simply rotates. The linear part rotates \emph{and} picks up $t\times R\omega$: the lever-arm effect. The same rule moves a twist between any two frames: $\xi_A=\Ad_{T_{AB}}\xi_B$.
\end{deriv}
\begin{example}{what does the camera feel when the robot turns?}
The base drives forward at $1$\,m/s while turning left at $0.5$\,rad/s: $\xi_B=(1,0,0;\ 0,0,0.5)$ (body twist). The forward-looking optical camera is mounted $0.3$\,m ahead and $0.5$\,m up ($T_{BC}$ with $R_{BC}^0$ from \S\ref{sec:rot:orientation}). Since $T_{WC}=T_{WB}T_{BC}$ with $T_{BC}$ constant, the camera's body twist is $\xi_C=\Ad_{T_{CB}}\xi_B$, which gives
\[
\xi_C=(-0.15,\ 0,\ 1.0;\ \ 0,\ -0.5,\ 0).
\]
Reading it in camera axes: moving forward along the optical axis at $1$\,m/s, sliding \emph{left} in the image at $0.15$\,m/s (lever arm $0.3$\,m $\times\,0.5$\,rad/s, because the camera is ahead of the turning centre), and rotating about $-y_C$ (image ``up'') at $0.5$\,rad/s. A visual-odometry or motion-blur model that ignored the lever arm would attribute that sideways slip to a nonexistent lateral velocity.
\end{example}

\subsection{Forward kinematics by the product of exponentials}\label{sec:rot:forward-kinematics-product}
\begin{deriv}{$T(\theta)=e^{[S_1]\theta_1}\cdots e^{[S_n]\theta_n}M$}
Let $M$ be the end-effector pose with all joints at zero, and $S_i$ the screw axis of joint $i$ expressed in the world frame \emph{at that home configuration}. Move only the last joint: the tool undergoes a screw about $S_n$, so $T=e^{[S_n]\theta_n}M$. Now move joint $n-1$: it carries joint $n$ and the tool rigidly, and its own axis has not been moved by joint $n$ (joint $n$ is downstream). So $T=e^{[S_{n-1}]\theta_{n-1}}e^{[S_n]\theta_n}M$. Induct to the base. No Denavit--Hartenberg frames are needed: only $n$ screw axes and one home pose.
\end{deriv}
\begin{figure}[h]
\centering
\begin{minipage}{0.36\linewidth}\includegraphics[width=\linewidth]{arm.pdf}\end{minipage}\hfill
\begin{minipage}{0.6\linewidth}
\caption{The 3-DOF elbow arm of the example at home (grey) and at $\theta=(30^\circ,45^\circ,-60^\circ)$ (blue). Each joint is described only by its screw axis at home; the product of exponentials moves the tool.}
\label{fig:rot:arm}
\end{minipage}
\end{figure}
\begin{example}{a 3-DOF elbow arm}
Link lengths $L_0=0.3$ (base height), $L_1=0.4$, $L_2=0.3$\,m; at home the arm points along $+x$. Joint 1: $\omega_1=e_z$ through the origin. Joints 2, 3: $\omega=-e_y$ (positive angle lifts the arm) through $q_2=(0,0,L_0)$ and $q_3=(L_1,0,L_0)$. With $v_i=-\omega_i\times q_i$:
\[
S_1=(0,0,0;\,0,0,1),\ S_2=(0.3,0,0;\,0,-1,0),\ S_3=(0.3,0,-0.4;\,0,-1,0),
\]
$M=\bmat{I&(0.7,0,0.3)\T\\0&1}$. For $\theta=(30^\circ,45^\circ,-60^\circ)$ the product gives tool position $(0.496,\,0.286,\,0.505)$\,m and tool $z$-axis $(0.224,0.129,0.966)$. Elementary check: horizontal reach $L_1\cos45^\circ+L_2\cos(-15^\circ)=0.573$, height $L_0+L_1\sin45^\circ+L_2\sin(-15^\circ)=0.505$, and $0.573(\cos30^\circ,\sin30^\circ)=(0.496,0.286)$.
\end{example}

\textbf{Velocity kinematics comes for free.} The spatial Jacobian has columns $J_1=S_1$ and $J_i=\Ad_{e^{[S_1]\theta_1}\cdots e^{[S_{i-1}]\theta_{i-1}}}S_i$: each joint's screw axis, carried to its \emph{current} location by the joints before it. For the arm above, $J_2=(0.260,0.150,0;\ 0.5,-0.866,0)$, $J_3=(0.505,0.291,-0.283;\ 0.5,-0.866,0)$. For $\dot\theta=(0.1,-0.2,0.3)$\,rad/s, $\xi_s=J\dot\theta$ and the tool-tip velocity $v_s+\omega_s\times p$ is $(0.0271,\,0.0818,\,-0.0276)$\,m/s, which matches finite differences of the forward kinematics.

\subsection{Pose optimisation: the $\SE$ version of ``tangent update, group state''}\label{sec:rot:pose-optimisation-version}
Exactly as in \S\ref{sec:rot:perturb}, estimators keep $T\in\SE$ and solve for a 6-vector correction $\delta=(\rho,\varphi)$, updating $T\leftarrow T\,\Exp(\delta)$. Using $\Exp(\delta)p\approx p+\rho+\varphi\times p$:
\[
\frac{\partial\,\big(T\Exp(\delta)\,p\big)}{\partial\delta}\Big|_{\delta=0}=R\,\big[\ I\ \ \big|\ -\skw{p}\ \big]\in\RR^{3\times6}.
\]
Chain this with the pinhole projection derivative and you have the analytic Jacobian of a reprojection error (bundle adjustment, PnP refinement); with a point-to-plane residual you have ICP. Twists, the adjoint and these Jacobians are the whole mathematical toolkit of modern SLAM, VIO, and robot calibration.
